%!TEX root = ./main.tex

\section[Act III: Security]{Act III: Antivirus on a Tower, Attacks without a Tower}

% SECTION TIMING: 43:00--72:00 (10 frames). Three papers + one teaser.
%   Paper E  43:00--53:00  5G-Spector (NDSS'24)      -- defenders enter through openness
%   Paper F  53:00--63:00  SigOver     (USENIX'19)   -- attack without a fake base station
%   Paper G  63:00--70:00  5G fronthaul (USENIX'24)  -- attack without any radio at all
%   Teaser   70:00--72:00  Sni5Gect    (USENIX'25)   -- ONE frame, then stop
% The act must end at 72:00; the wrap-up owns 72:00--75:00.
%
% FRAMING BOUNDARY, enforced throughout: everything here is published, peer-reviewed
% research described at the CONCEPT level. No commands, configs, or step-by-step
% recipes. Nothing here claims that anyone in the room is being attacked. If a student
% asks "can I do this", the answer is that the papers exist, the code is the authors'
% to release, and this course teaches how the architecture makes attacks possible so
% that defenders understand them.
%
% LAYOUT RULE: every bullet fits on one line. Shorten the text, do not let it wrap.

% ===========================================================================
% PAPER E -- 5G-Spector (NDSS 2024)
% ===========================================================================

% Teaching note (2 min): callback to Class 4. PCs were not born with antivirus; they
% got it once they became programmable platforms anyone could ship software to. Ask
% the \ask out loud BEFORE advancing -- most students will say "no, a tower is not a
% computer". That is the trap the next slide springs: O-RAN made it one.
\begin{frame}{Can we install antivirus on a cell tower?}
  \vspace{0.4em}
  \begin{columns}[T,onlytextwidth]
    \column{0.48\textwidth}
      \large\textbf{Why did PCs get antivirus?}

      \vspace{0.5em}
      \normalsize
      Not because they were computers.

      \vspace{0.4em}
      Because they became \textbf{programmable platforms} anyone could ship
      software to --- including defenders.
    \column{0.48\textwidth}
      \large\textbf{And the RAN?}

      \vspace{0.5em}
      \normalsize
      For decades a base station ran one vendor's locked firmware.

      \vspace{0.4em}
      Last class you learned O-RAN opened it: \textbf{RIC + xApps}.
  \end{columns}

  \vspace{0.8em}
  \ask{If the tower is now a platform, can you install a security app on it?}
\end{frame}

% Teaching note (3 min): the correction that matters -- this is NOT "AI catches
% attacks". It is telemetry + rules, the same shape as a classic intrusion detector.
% MOBIFLOW = a stream of fine-grained RAN events (say "NetFlow for the RAN" and let it
% land). MOBIEXPERT = an xApp on the RIC that runs production rules over that stream.
% Draw the pipe left to right; the xApp box is the one they drew themselves last class.
\begin{frame}{Not deep learning --- telemetry plus rules}
  \vspace{0.3em}
  \centering
  \begin{tikzpicture}[scale=0.9,transform shape]
    \node[netnode,minimum width=1.7cm] (ran) at (0,0) {RAN\\events};
    \node[oranode,minimum width=2.2cm] (flow) at (3.0,0) {MobiFlow\\telemetry};
    \node[attacknode,minimum width=2.3cm] (exp) at (6.6,0) {MobiExpert\\rule xApp};
    \node[corenode,minimum width=1.7cm] (alert) at (9.8,0) {alert};
    \draw[flow] (ran) -- (flow);
    \draw[flow] (flow) -- (exp);
    \draw[flow] (exp) -- (alert);
    \node[font=\scriptsize,text=softgray,anchor=north] at (3.0,-0.6)
      {``NetFlow for the RAN''};
    \node[font=\scriptsize,text=softgray,anchor=north] at (6.6,-0.6)
      {runs on the RIC};
    \node[font=\tiny,text=softgray,anchor=north] at (4.9,-1.05)
      {NetFlow = the per-connection logs network admins already collect};
  \end{tikzpicture}

  \vspace{0.7em}
  \small
  A stream of fine-grained cellular telemetry feeds a rule engine that watches
  for known layer-3 \expand{the signaling layer: attach, incoming-call alerts,
  handover messages} attack patterns. No model to train, no black box.

  \vspace{0.6em}
  \papercard{5G-Spector: An O-RAN Compliant Layer-3 Cellular Attack Detection Service}
    {Haohuang Wen, Zhiqiang Lin (Ohio State); Porras, Yegneswaran, Gehani (SRI)}
    {NDSS 2024}
    {https://www.ndss-symposium.org/ndss-paper/5g-spector-an-o-ran-compliant-layer-3-cellular-attack-detection-service/}
\end{frame}

% Teaching note (3 min): land the numbers, then the payoff. The point of "extensible to
% 11 new variants" is that operators write new rules without redeploying the network --
% the platform property they learned about, doing real work. The award line is not
% bragging; it means the artifact was reproducible enough for a committee to run.
\begin{frame}{What it caught, and what it cost}
  \vspace{0.4em}
  \small
  \begin{itemize}\itemsep6pt
    % SOURCE: https://www.ndss-symposium.org/ndss-paper/5g-spector-an-o-ran-compliant-layer-3-cellular-attack-detection-service/
    \item Detects \textbf{7 classes} of layer-3 cellular attacks in real time.
    % SOURCE: https://www.ndss-symposium.org/ndss-paper/5g-spector-an-o-ran-compliant-layer-3-cellular-attack-detection-service/
    \item Shown to scale to \textbf{11 previously unknown attacks} by writing new rules.
    % SOURCE: https://www.ndss-symposium.org/ndss-paper/5g-spector-an-o-ran-compliant-layer-3-cellular-attack-detection-service/
    \item CPU overhead under \textbf{2\%}; memory under \textbf{100\,MB}.
    % SOURCE: https://www.5gsec.com/post/award-5g-spector-won-the-ndss-24-distinguished-artifact-award
    % SOURCE: https://www.ndss-symposium.org/news/ndss-2024-largest-ndss-signals-continued-importance-of-distributed-systems-security/
    \item Won an \textbf{NDSS 2024 Distinguished Artifact Award}.
  \end{itemize}

  \vspace{0.6em}
  \qa{So was the xApp box on our O-RAN slide just a PowerPoint concept?}
     {No. Someone shipped a working security function inside it --- and a
      committee reran it. The box on your diagram is a place code lives.}
\end{frame}

% Teaching note (say before leaving this frame): 5G-Spector's antivirus trusts
% the telemetry the network feeds it. Pocket that for twenty minutes.

% ===========================================================================
% PAPER F -- SigOver (USENIX Security 2019)
% ===========================================================================

% Teaching note (2 min): the pivot. 5G-Spector detects attacks -- but what does a modern
% attack even look like? Everyone in the room "knows" from Class 4 that you attack
% cellular by standing up a fake base station. Say the assumption aloud, then break it.
\begin{frame}{Attacking without building a tower}
  \vspace{0.5em}
  % Wording of the assumption below must stay IDENTICAL to item 5 on the
  % "Five things the diagram never says out loud" frame in 01_assumptions.tex.
  \brokenassumption{Attacking the network requires \textbf{standing up} a fake base station.}
    {Whispering in perfect sync with the real tower, slightly louder.}

  \vspace{0.7em}
  \normalsize
  A rogue base station is loud, obvious, and easy to look for --- it is a
  whole tower that should not be there.

  \vspace{0.6em}
  SigOver never builds one. It \textbf{synchronizes} to the legitimate LTE
  downlink and overwrites only the subframes it chooses
  \expand{slices of the LTE frame, a millisecond each}.

  \vspace{0.5em}
  {\small\textcolor{softgray}{Think: two voices, identical timing --- next
   slide.}\par}

  \vspace{0.5em}
  \takeaway{You do not need a new tower to lie ---\\you need better timing.}

  \glossline{LTE = the 4G radio standard \quad downlink = the tower-to-phone direction}
\end{frame}

% Teaching note (3 min): the analogy is the whole slide. Two voices, IDENTICAL timing,
% ear keeps the louder. Then the three-way distinction -- read one line each and let the
% contrast do the teaching. The honesty box is mandatory: this is LTE, 2019, and it
% targets BROADCAST messages, which are never integrity protected -- before or after
% authentication.
\begin{frame}{Jamming, spoofing, overshadowing}
  \vspace{0.3em}
  \small
  \renewcommand{\arraystretch}{1.4}
  \begin{tabular}{@{}p{0.26\textwidth}p{0.66\textwidth}@{}}
    \toprule
    \textbf{Jamming} & You hear \textbf{nothing} --- the channel is drowned out. \\
    \textbf{Spoofing (fake BS)} & You \textbf{switch towers} --- a new cell lures you over. \\
    \textbf{\textcolor{attackred}{Overshadowing}} & You hear the \textbf{attacker} but believe it is your tower. \\
    \bottomrule
  \end{tabular}

  \vspace{0.8em}
  \normalsize
  Two people speak to you with \textbf{identical timing}. Your ear keeps the
  louder one --- and never notices there were two.

  \vspace{0.6em}
  {\scriptsize\textcolor{softgray}{Scope: this is LTE (2019). It targets LTE
   \textbf{broadcast} messages, which are never integrity protected --- before
   or after authentication. It is not a claim about 5G user traffic.}\par}
\end{frame}

% Teaching note (2 min): numbers, then play the demo. The 98%-vs-80% comparison is the
% punchline: overshadowing at a 3 dB edge beats a fake base station that is 35 dB louder,
% because the victim's receiver was built to lock onto the best-timed signal, not the
% loudest stranger. Play ~90 s of the talk demo, then move on.
\begin{frame}{Hiding in plain signal}
  \vspace{0.2em}
  \small
  \begin{itemize}\itemsep5pt
    % SOURCE: https://www.usenix.org/conference/usenixsecurity19/presentation/yang-hojoon
    \item A signal only \textbf{3\,dB} \expand{about twice the power} stronger wins.
    % SOURCE: https://www.usenix.org/conference/usenixsecurity19/presentation/yang-hojoon
    \item Success rate \textbf{98\%} --- vs \textbf{80\%} for a fake base station.
    % SOURCE: https://www.usenix.org/conference/usenixsecurity19/presentation/yang-hojoon
    \item That fake base station needed a \textbf{35\,dB} \expand{about 3000\(\times\)} edge.
    \item Timing is why 3\,dB is enough: a time-aligned signal wins the receiver's capture.
    \item A fake tower has no such luck: it must out-shout cell selection.
  \end{itemize}

  \glossline{capture = the radio locks onto the strongest aligned signal and
    treats the rest as noise\\
    cell selection = how your phone picks which tower to camp on}

  \papercard{Hiding in Plain Signal: Physical Signal Overshadowing Attack on LTE}
    {Hojoon Yang, Sangwook Bae, Mincheol Son, Hongil Kim, Song Min Kim, Yongdae Kim (KAIST)}
    {USENIX Security 2019}
    {https://www.usenix.org/conference/usenixsecurity19/presentation/yang-hojoon}

  \watch{https://www.youtube.com/watch?v=xvIaJTJFRJQ}{USENIX Sec'19 talk --- play the demo, \(\sim\)90\,s}
\end{frame}

% ===========================================================================
% PAPER G -- 5G Fronthaul Integrity (USENIX Security 2024)
% ===========================================================================

% CUT FIRST if running long: this frame's content can be delivered verbally over
% the next one.
% Teaching note (2 min): the reversal. SigOver was intensely physical -- SDR, timing,
% power in dB. Now flip it: what if the attacker sends no radio at all? Let the question
% sit before advancing. It reframes "attack" from the air interface to the wiring.
\begin{frame}{What if the attacker never transmits a radio signal?}
  \vspace{0.8em}
  \large
  SigOver was intensely physical: a radio, microsecond timing, power measured
  in decibels.

  \vspace{1.0em}
  \normalsize
  Now turn the question inside out.

  \vspace{0.8em}
  \ask{What could you break from software, without touching the air at all?}

  \vspace{0.6em}
  \small
  \textcolor{softgray}{The answer is on a link you drew last class and never
  thought twice about.}
\end{frame}

% Teaching note (3 min): recall THEIR diagram. Before disaggregation, fronthaul was
% dedicated point-to-point fiber (CPRI) -- effectively private wiring, physically
% trusted. Disaggregated 5G: it is eCPRI over switched, packet-based Ethernet. To
% hold microsecond latency budgets, integrity protection is not mandatory in the
% fronthaul standards -- careful framing: this is about what the standards require
% and what deployments do, not a claim every network is wide open.
\begin{frame}{The link inside the tower is now a network}
  \vspace{0.2em}
  \centering
  \begin{tikzpicture}[scale=0.9,transform shape]
    \node[netnode,minimum width=1.5cm] (ru) at (0,0) {RU};
    \node[netnode,minimum width=1.5cm] (du) at (4.2,0) {DU};
    \draw[flow] (ru) -- (du);
    \node[font=\scriptsize,text=softgray,anchor=south] at (2.1,0.15)
      {eCPRI over Ethernet};
    \node[attacknode,minimum width=2.0cm] (atk) at (2.1,-1.5) {software\\attacker};
    \draw[flow,draw=attackred,dashed] (atk) -- (2.1,-0.15);
  \end{tikzpicture}

  \vspace{0.4em}
  \small
  \begin{columns}[T,onlytextwidth]
    \column{0.47\textwidth}
      \textbf{Before:} fronthaul was dedicated point-to-point fiber (CPRI) ---
      effectively private wiring, physically trusted.
    \column{0.47\textwidth}
      \textbf{Disaggregated 5G:} switched, packet-based Ethernet (eCPRI) --- and
      to stay fast enough, nobody is required to check packets are genuine.
  \end{columns}

  \vspace{0.3em}
  {\scriptsize\textcolor{softgray}{Framing: the fronthaul standards do not
   require integrity protection, citing performance overhead and assumed low
   risk. That assumption is what the paper tests.}\par}

  \glossline{RU = Radio Unit (antenna end) \quad DU = Distributed Unit (baseband
    computer) \quad eCPRI = the packet protocol between them}

  \vspace{0.2em}
  \brokenassumption{The links inside the base station are trustworthy}
    {a software attacker on the fronthaul Ethernet}
\end{frame}

% Teaching note (3 min): impact, then the tension that closes the act. Numbers are from
% a commercial-grade multi-cell testbed with just two compromised cells and four users --
% say that scale out loud so nobody hears "the whole country". Then the closing Q/A ties
% Act III together: 5G-Spector said openness lets defenders in; this says disaggregation
% lets attackers in. Both true. That is why it is research, not a settled answer.
\begin{frame}{What software alone could do}
  \small
  % SOURCE: https://www.usenix.org/conference/usenixsecurity24/presentation/xing-jiarong
  On a commercial-grade, multi-cell testbed --- two compromised cells, four users:
  \begin{itemize}\itemsep2pt
    % SOURCE: https://www.usenix.org/conference/usenixsecurity24/presentation/xing-jiarong
    \item User performance degraded by more than \textbf{80\%}.
    % SOURCE: https://www.usenix.org/conference/usenixsecurity24/presentation/xing-jiarong
    \item Targeted users \textbf{prevented from ever attaching} to the cell.
    % SOURCE: https://www.usenix.org/conference/usenixsecurity24/presentation/xing-jiarong
    \item Signaling storms \expand{floods of junk control messages} over \textbf{2500 msgs/min}.
  \end{itemize}

  \vspace{0.1em}
  \papercard{On the Criticality of Integrity Protection in 5G Fronthaul Networks}
    {Jiarong Xing (Rice); Yoo (Princeton); Foukas (Microsoft); Kim (UT Austin); Reiter (Duke)}
    {USENIX Security 2024}
    {https://www.usenix.org/conference/usenixsecurity24/presentation/xing-jiarong}

  \qa{The same opened seams let 5G-Spector in --- and this attacker in. Safer, or
      more dangerous?}
     {Both. Openness is a trade, not a verdict --- which is why this is research,
      not settled fact.}
\end{frame}

% Teaching note (say while leaving this frame): the fronthaul attacker sits on the
% wires that telemetry rides. Could it lie to the antivirus?
% (The \watch talk link was cut from this frame to fix an overfull page; the talk
% is one click away from the papercard's USENIX page if you want it.)

% ===========================================================================
% TEASER -- Sni5Gect (USENIX Security 2025) -- ONE frame, then stop.
% ===========================================================================

% Teaching note (2 min): do NOT explain the mechanism. This is a hand-off to the paper
% practicum -- one of the students will present it. The single idea to leave in the room:
% SigOver's lineage did not die with LTE. The same "no rogue tower" move reappears on 5G
% NR. Name the numbers, name the promise, stop.
\begin{frame}{The lineage did not die with LTE}
  \normalsize
  SigOver's core idea --- attack without a rogue tower --- came back for 5G NR.

  \vspace{0.4em}
  \small
  % SOURCE: https://www.usenix.org/conference/usenixsecurity25/presentation/luo-shijie
  \begin{itemize}\itemsep4pt
    \item Real-time sniffing of \textbf{pre-authentication 5G NR} traffic.
    \item Targeted downlink message injection, \textbf{no rogue base station}.
    % SOURCE: https://www.usenix.org/conference/usenixsecurity25/presentation/luo-shijie
    \item \textbf{70--90\%} injection success at up to \textbf{\(\sim\)20\,m}.
  \end{itemize}

  \glossline{5G NR = the 5G radio standard \quad pre-authentication = messages
    sent before the phone proves its identity}

  \vspace{0.1em}
  % SOURCE: https://zenodo.org/records/15601773 (Zenodo artifact confirms all
  % four authors are SUTD)
  \papercard{Sni5Gect: A Practical Approach to Inject aNRchy into 5G NR}
    {Shijie Luo, Matheus E. Garbelini, Sudipta Chattopadhyay, Jianying Zhou (SUTD)}
    {USENIX Security 2025}
    {https://www.usenix.org/conference/usenixsecurity25/presentation/luo-shijie}

  \vspace{0.1em}
  \takeaway{You will meet this paper again ---\\one of you presents it in the practicum.}
\end{frame}
